% Options for packages loaded elsewhere
\PassOptionsToPackage{unicode}{hyperref}
\PassOptionsToPackage{hyphens}{url}
\documentclass[
  11pt,
  a4paper,
]{article}
\usepackage{xcolor}
\usepackage[margin=2.5cm]{geometry}
\usepackage{amsmath,amssymb}
\setcounter{secnumdepth}{-\maxdimen} % remove section numbering
\usepackage{iftex}
\ifPDFTeX
  \usepackage[T1]{fontenc}
  \usepackage[utf8]{inputenc}
  \usepackage{textcomp} % provide euro and other symbols
\else % if luatex or xetex
  \usepackage{unicode-math} % this also loads fontspec
  \defaultfontfeatures{Scale=MatchLowercase}
  \defaultfontfeatures[\rmfamily]{Ligatures=TeX,Scale=1}
\fi
\usepackage{lmodern}
\ifPDFTeX\else
  % xetex/luatex font selection
\fi
% Use upquote if available, for straight quotes in verbatim environments
\IfFileExists{upquote.sty}{\usepackage{upquote}}{}
\IfFileExists{microtype.sty}{% use microtype if available
  \usepackage[]{microtype}
  \UseMicrotypeSet[protrusion]{basicmath} % disable protrusion for tt fonts
}{}
\makeatletter
\@ifundefined{KOMAClassName}{% if non-KOMA class
  \IfFileExists{parskip.sty}{%
    \usepackage{parskip}
  }{% else
    \setlength{\parindent}{0pt}
    \setlength{\parskip}{6pt plus 2pt minus 1pt}}
}{% if KOMA class
  \KOMAoptions{parskip=half}}
\makeatother
\usepackage{color}
\usepackage{fancyvrb}
\newcommand{\VerbBar}{|}
\newcommand{\VERB}{\Verb[commandchars=\\\{\}]}
\DefineVerbatimEnvironment{Highlighting}{Verbatim}{commandchars=\\\{\}}
% Add ',fontsize=\small' for more characters per line
\usepackage{framed}
\definecolor{shadecolor}{RGB}{248,248,248}
\newenvironment{Shaded}{\begin{snugshade}}{\end{snugshade}}
\newcommand{\AlertTok}[1]{\textcolor[rgb]{0.94,0.16,0.16}{#1}}
\newcommand{\AnnotationTok}[1]{\textcolor[rgb]{0.56,0.35,0.01}{\textbf{\textit{#1}}}}
\newcommand{\AttributeTok}[1]{\textcolor[rgb]{0.13,0.29,0.53}{#1}}
\newcommand{\BaseNTok}[1]{\textcolor[rgb]{0.00,0.00,0.81}{#1}}
\newcommand{\BuiltInTok}[1]{#1}
\newcommand{\CharTok}[1]{\textcolor[rgb]{0.31,0.60,0.02}{#1}}
\newcommand{\CommentTok}[1]{\textcolor[rgb]{0.56,0.35,0.01}{\textit{#1}}}
\newcommand{\CommentVarTok}[1]{\textcolor[rgb]{0.56,0.35,0.01}{\textbf{\textit{#1}}}}
\newcommand{\ConstantTok}[1]{\textcolor[rgb]{0.56,0.35,0.01}{#1}}
\newcommand{\ControlFlowTok}[1]{\textcolor[rgb]{0.13,0.29,0.53}{\textbf{#1}}}
\newcommand{\DataTypeTok}[1]{\textcolor[rgb]{0.13,0.29,0.53}{#1}}
\newcommand{\DecValTok}[1]{\textcolor[rgb]{0.00,0.00,0.81}{#1}}
\newcommand{\DocumentationTok}[1]{\textcolor[rgb]{0.56,0.35,0.01}{\textbf{\textit{#1}}}}
\newcommand{\ErrorTok}[1]{\textcolor[rgb]{0.64,0.00,0.00}{\textbf{#1}}}
\newcommand{\ExtensionTok}[1]{#1}
\newcommand{\FloatTok}[1]{\textcolor[rgb]{0.00,0.00,0.81}{#1}}
\newcommand{\FunctionTok}[1]{\textcolor[rgb]{0.13,0.29,0.53}{\textbf{#1}}}
\newcommand{\ImportTok}[1]{#1}
\newcommand{\InformationTok}[1]{\textcolor[rgb]{0.56,0.35,0.01}{\textbf{\textit{#1}}}}
\newcommand{\KeywordTok}[1]{\textcolor[rgb]{0.13,0.29,0.53}{\textbf{#1}}}
\newcommand{\NormalTok}[1]{#1}
\newcommand{\OperatorTok}[1]{\textcolor[rgb]{0.81,0.36,0.00}{\textbf{#1}}}
\newcommand{\OtherTok}[1]{\textcolor[rgb]{0.56,0.35,0.01}{#1}}
\newcommand{\PreprocessorTok}[1]{\textcolor[rgb]{0.56,0.35,0.01}{\textit{#1}}}
\newcommand{\RegionMarkerTok}[1]{#1}
\newcommand{\SpecialCharTok}[1]{\textcolor[rgb]{0.81,0.36,0.00}{\textbf{#1}}}
\newcommand{\SpecialStringTok}[1]{\textcolor[rgb]{0.31,0.60,0.02}{#1}}
\newcommand{\StringTok}[1]{\textcolor[rgb]{0.31,0.60,0.02}{#1}}
\newcommand{\VariableTok}[1]{\textcolor[rgb]{0.00,0.00,0.00}{#1}}
\newcommand{\VerbatimStringTok}[1]{\textcolor[rgb]{0.31,0.60,0.02}{#1}}
\newcommand{\WarningTok}[1]{\textcolor[rgb]{0.56,0.35,0.01}{\textbf{\textit{#1}}}}
\usepackage{longtable,booktabs,array}
\usepackage{calc} % for calculating minipage widths
% Correct order of tables after \paragraph or \subparagraph
\usepackage{etoolbox}
\makeatletter
\patchcmd\longtable{\par}{\if@noskipsec\mbox{}\fi\par}{}{}
\makeatother
% Allow footnotes in longtable head/foot
\IfFileExists{footnotehyper.sty}{\usepackage{footnotehyper}}{\usepackage{footnote}}
\makesavenoteenv{longtable}
\ifLuaTeX
\usepackage[bidi=basic]{babel}
\else
\usepackage[bidi=default]{babel}
\fi
\babelprovide[main,import]{spanish}
% get rid of language-specific shorthands (see #6817):
\let\LanguageShortHands\languageshorthands
\def\languageshorthands#1{}
\setlength{\emergencystretch}{3em} % prevent overfull lines
\providecommand{\tightlist}{%
  \setlength{\itemsep}{0pt}\setlength{\parskip}{0pt}}
\usepackage{array}
\usepackage{etoolbox}
\usepackage{booktabs}
\setlength{\tabcolsep}{4pt}
\BeforeBeginEnvironment{longtable}{\footnotesize}
\BeforeBeginEnvironment{tabular}{\footnotesize}
\renewcommand{\arraystretch}{1.1}
\usepackage{bookmark}
\IfFileExists{xurl.sty}{\usepackage{xurl}}{} % add URL line breaks if available
\urlstyle{same}
\hypersetup{
  pdflang={es},
  hidelinks,
  pdfcreator={LaTeX via pandoc}}

\author{}
\date{}

\begin{document}

{
\setcounter{tocdepth}{3}
\tableofcontents
}
\subsection{INTRODUCCIÓN}\label{introducciuxf3n}

Python se ha consolidado como uno de los lenguajes de programación más
relevantes en el campo de la seguridad informática. Su sintaxis clara,
la disponibilidad de miles de librerías especializadas y su extensa
comunidad lo convierten en la herramienta de elección para analistas,
investigadores y profesionales de la ciberseguridad defensiva en todo el
mundo.

El presente trabajo de investigación monográfico examina tres librerías
fundamentales del ecosistema Python orientadas a la ciberseguridad:
yara-python, hashlib e impacket. Estas herramientas permiten,
respectivamente, detectar y clasificar malware mediante reglas de
reconocimiento de patrones, verificar la integridad de archivos y
mensajes mediante funciones hash criptográficas, y analizar protocolos
de red como SMB, Kerberos y LDAP en entornos de laboratorio.

Como componente práctico principal, se desarrolló un caso real de
análisis forense sobre dos aplicaciones Android (APK): una original
(CairosOriginal.apk) y otra infectada con un payload de Metasploit
(CairosVenom.apk). Se utilizó hashlib para calcular y comparar los
hashes criptográficos de ambos archivos, evidenciando el efecto
avalancha, y se implementó un detector basado en yara-python y
Androguard capaz de identificar automáticamente el APK envenenado
mediante el análisis de permisos, subpaquetes inyectados y firmas
DEX. Todo el código fuente, los APK de prueba y este documento
se encuentran disponibles en el repositorio del proyecto.

El informe se organiza siguiendo el esquema de investigación
institucional: inicia con el planteamiento del problema y los objetivos,
continúa con el marco teórico y el desarrollo detallado de cada
librería, incluye una práctica ejecutada en Google Colab, describe un
caso práctico en contexto real y concluye con el análisis comparativo y
las conclusiones del grupo.

Este trabajo busca que los estudiantes no solo conozcan la definición de
las librerías, sino que comprendan su utilidad práctica dentro de
operaciones de ciberseguridad defensiva, siempre respetando marcos
éticos y legales vigentes.

\subsubsection{Propósito de la guía}\label{propuxf3sito-de-la-guuxeda}

\begin{itemize}
\tightlist
\item
  Orientar la elaboración del trabajo monográfico sobre librerías de
  Python aplicadas a la ciberseguridad.
\item
  Relacionar la investigación teórica con una práctica funcional
  ejecutada en notebook.
\item
  Promover el análisis de casos reales y la interpretación de resultados
  en contextos de seguridad informática.
\item
  Demostrar el valor defensivo de yara-python, hashlib e impacket en
  entornos académicos y profesionales.
\end{itemize}

\subsection{CAPÍTULO I. PLANTEAMIENTO GENERAL DEL
TEMA}\label{capuxedtulo-i.-planteamiento-general-del-tema}

\subsubsection{1.1. Descripción del tema}\label{descripciuxf3n-del-tema}

Las librerías de Python son componentes de software precompilados que
ofrecen funcionalidades específicas listas para ser integradas en
proyectos de programación. En el ámbito de la ciberseguridad, estas
librerías reducen drásticamente el tiempo necesario para implementar
mecanismos defensivos, ya que los analistas pueden enfocarse en la
lógica de detección y respuesta en lugar de construir desde cero cada
función criptográfica o de análisis de protocolos.

El problema que este trabajo aborda es la necesidad de contar con
herramientas de detección de amenazas, verificación de integridad y
análisis de tráfico de red que sean accesibles, documentadas y
aplicables en entornos educativos. Las tres librerías seleccionadas ---
yara-python, hashlib e impacket --- responden directamente a esta
necesidad: yara-python permite identificar malware mediante firmas de
patrones; hashlib garantiza la integridad de archivos mediante funciones
hash como SHA-256 y MD5; e impacket facilita el análisis y la simulación
de protocolos de red como SMB y Kerberos en laboratorio.

\subsubsection{1.2. Importancia del uso de librerías en
Python}\label{importancia-del-uso-de-libreruxedas-en-python}

Las librerías permiten reutilizar código ampliamente probado por la
comunidad, reducir errores humanos en implementaciones críticas de
seguridad y acceder a algoritmos complejos con interfaces de
programación simples. En ciberseguridad, donde un error de
implementación puede comprometer la totalidad de un sistema, el uso de
librerías estandarizadas resulta indispensable.

Además, las librerías facilitan el trabajo colaborativo entre equipos de
seguridad, la integración con plataformas de análisis de amenazas y la
automatización de tareas repetitivas como el cálculo masivo de hashes o
el escaneo de archivos con reglas YARA.

\subsubsection{1.3. Área de aplicación
seleccionada}\label{uxe1rea-de-aplicaciuxf3n-seleccionada}

Este trabajo se ubica en el área de ciberseguridad defensiva. Las
librerías seleccionadas contribuyen de forma complementaria a la
protección de sistemas: hashlib para integridad, yara-python para
detección de código malicioso e impacket para análisis de protocolos de
red en entornos controlados.

\begin{longtable}[]{@{}
  >{\raggedright\arraybackslash}p{(\linewidth - 2\tabcolsep) * \real{0.5000}}
  >{\raggedright\arraybackslash}p{(\linewidth - 2\tabcolsep) * \real{0.5000}}@{}}
\toprule\noalign{}
\begin{minipage}[b]{\linewidth}\raggedright
Pregunta guía
\end{minipage} & \begin{minipage}[b]{\linewidth}\raggedright
Respuesta desarrollada
\end{minipage} \\
\midrule\noalign{}
\endhead
\bottomrule\noalign{}
\endlastfoot
¿Qué área eligieron? & Ciberseguridad defensiva: yara-python, hashlib e
impacket. \\
¿Por qué eligieron esa área? & La ciberseguridad es una competencia
crítica en la formación tecnológica actual. Las amenazas digitales
crecen constantemente y las organizaciones requieren profesionales
capacitados en herramientas defensivas. \\
¿Qué problema real puede resolver? & Detección de malware en archivos
descargados, verificación de integridad de software distribuido en red,
y análisis de tráfico SMB/Kerberos en auditorías internas de
seguridad. \\
¿Qué resultados esperan obtener? & Generación de hashes verificables,
detección de patrones maliciosos en archivos de prueba, y captura de
paquetes de protocolos de red en laboratorio. \\
\end{longtable}

\subsection{CAPÍTULO II. OBJETIVOS DEL
TRABAJO}\label{capuxedtulo-ii.-objetivos-del-trabajo}

\subsubsection{2.1. Objetivo general}\label{objetivo-general}

Analizar el uso, características y aplicación práctica de las librerías
yara-python, hashlib e impacket de Python orientadas a la ciberseguridad
defensiva, mediante una investigación monográfica y el desarrollo de un
notebook funcional en Google Colab.

\subsubsection{2.2. Objetivos
específicos}\label{objetivos-especuxedficos}

\begin{enumerate}
\def\labelenumi{\arabic{enumi}.}
\tightlist
\item
  Identificar las principales características, funciones y métodos de
  las librerías yara-python, hashlib e impacket.
\item
  Explicar el funcionamiento básico, la instalación y la importación de
  cada librería, junto con sus métodos principales.
\item
  Comparar las tres librerías según su uso principal, nivel de
  dificultad, ventajas, limitaciones y campo de aplicación real.
\item
  Desarrollar una práctica en Google Colab utilizando las tres librerías
  para resolver un problema de ciberseguridad defensiva.
\item
  Interpretar los resultados obtenidos y relacionarlos con un caso
  práctico en contexto de seguridad informática institucional.
\end{enumerate}

\subsubsection{2.3. Coherencia de
objetivos}\label{coherencia-de-objetivos}

\begin{longtable}[]{@{}
  >{\raggedright\arraybackslash}p{(\linewidth - 2\tabcolsep) * \real{0.5000}}
  >{\raggedright\arraybackslash}p{(\linewidth - 2\tabcolsep) * \real{0.5000}}@{}}
\toprule\noalign{}
\begin{minipage}[b]{\linewidth}\raggedright
Elemento
\end{minipage} & \begin{minipage}[b]{\linewidth}\raggedright
Verificación
\end{minipage} \\
\midrule\noalign{}
\endhead
\bottomrule\noalign{}
\endlastfoot
Título & Menciona librerías de Python y el área de ciberseguridad. \\
Objetivo general & Coincide con la finalidad del trabajo: análisis e
implementación práctica. \\
Objetivos específicos & Organizan el desarrollo del informe de lo
teórico a lo práctico. \\
Parte práctica & Evidenciará el cumplimiento de todos los objetivos
mediante código ejecutado. \\
\end{longtable}

\subsection{CAPÍTULO III. MARCO
TEÓRICO}\label{capuxedtulo-iii.-marco-teuxf3rico}

\subsubsection{3.1. Concepto de Python}\label{concepto-de-python}

Python es un lenguaje de programación de alto nivel, interpretado, de
tipado dinámico y de propósito general, creado por Guido van Rossum y
publicado por primera vez en 1991. Se caracteriza por su sintaxis
legible que prioriza la claridad del código, lo que lo convierte en uno
de los lenguajes más utilizados tanto en educación como en entornos
profesionales. En la actualidad, Python lidera índices de popularidad
como el TIOBE Index y PyPL, y es la primera opción en campos como la
ciencia de datos, la inteligencia artificial, la automatización y la
ciberseguridad.

En el ámbito de la seguridad informática, Python permite desarrollar
herramientas de análisis de malware, scripts de automatización de
auditorías, sistemas de monitoreo de red y motores de cifrado, gracias a
su extensa biblioteca estándar y al ecosistema de librerías de terceros
disponibles en PyPI (Python Package Index).

\subsubsection{3.2. Concepto de librería en
programación}\label{concepto-de-libreruxeda-en-programaciuxf3n}

Una librería es un conjunto de módulos, clases, funciones y recursos
precompilados que permiten resolver tareas específicas dentro de un
programa sin necesidad de implementar cada funcionalidad desde cero. Las
librerías encapsulan lógica compleja y son distribuidas de forma que
cualquier desarrollador pueda integrarlas mediante una simple
instrucción de importación.

Ejemplo de importación de una librería en Python:

\begin{Shaded}
\begin{Highlighting}[]
\ImportTok{import}\NormalTok{ hashlib}
\CommentTok{\# La librería hashlib permite calcular funciones hash criptográficas}
\CommentTok{\# como SHA{-}256 o MD5.}
\end{Highlighting}
\end{Shaded}

\subsubsection{3.3. Clasificación de librerías de
Python}\label{clasificaciuxf3n-de-libreruxedas-de-python}

\begin{longtable}[]{@{}
  >{\raggedright\arraybackslash}p{(\linewidth - 4\tabcolsep) * \real{0.3333}}
  >{\raggedright\arraybackslash}p{(\linewidth - 4\tabcolsep) * \real{0.3333}}
  >{\raggedright\arraybackslash}p{(\linewidth - 4\tabcolsep) * \real{0.3333}}@{}}
\toprule\noalign{}
\begin{minipage}[b]{\linewidth}\raggedright
Clasificación
\end{minipage} & \begin{minipage}[b]{\linewidth}\raggedright
Descripción
\end{minipage} & \begin{minipage}[b]{\linewidth}\raggedright
Ejemplos
\end{minipage} \\
\midrule\noalign{}
\endhead
\bottomrule\noalign{}
\endlastfoot
Analítica de datos & Permiten limpiar, transformar, analizar y
visualizar información estructurada. & Pandas, NumPy, Matplotlib,
Seaborn \\
Computación gráfica & Permiten trabajar con imágenes, videojuegos,
gráficos 2D/3D o visualizaciones. & OpenCV, Pillow, Pygame, PyOpenGL \\
Inteligencia artificial & Permiten crear modelos que aprenden de datos y
realizan predicciones. & Scikit-learn, TensorFlow, PyTorch, Keras \\
Procesamiento de lenguaje natural & Permiten analizar textos, palabras,
opiniones y lenguaje humano. & NLTK, spaCy, Transformers \\
Ciberseguridad defensiva & Permiten cifrar, verificar integridad,
analizar archivos o protocolos de red. & hashlib, yara-python, impacket,
Cryptography, Scapy \\
\end{longtable}

\subsubsection{3.4. Importancia de las librerías en problemas
reales}\label{importancia-de-las-libreruxedas-en-problemas-reales}

En el contexto de la ciberseguridad, las librerías de Python tienen un
impacto directo en la capacidad de respuesta ante amenazas. hashlib
permite a una organización verificar si un archivo descargado ha sido
alterado comparando su hash con el valor original publicado por el
proveedor. yara-python permite a un analista de SOC (Security Operations
Center) escanear miles de archivos en segundos para detectar firmas de
familias de malware conocidas. impacket, por su parte, es utilizada en
auditorías de seguridad para simular ataques sobre protocolos de red y
detectar vulnerabilidades antes de que actores maliciosos las exploten.

\subsection{CAPÍTULO IV. DESARROLLO DE LAS LIBRERÍAS
SELECCIONADAS}\label{capuxedtulo-iv.-desarrollo-de-las-libreruxedas-seleccionadas}

\subsubsection{4.1. LIBRERÍA N.º 1:
hashlib}\label{libreruxeda-n.uxba-1-hashlib}

\paragraph{4.1.1. Nombre de la librería}\label{nombre-de-la-libreruxeda}

Nombre oficial: hashlib. Se importa directamente bajo el mismo nombre:
\texttt{import\ hashlib}. Es parte de la biblioteca estándar de Python,
por lo que no requiere instalación adicional.

\paragraph{4.1.2. Definición de la
librería}\label{definiciuxf3n-de-la-libreruxeda}

hashlib es una librería estándar de Python que proporciona interfaces
para los algoritmos de hash criptográficos más utilizados, incluyendo
MD5, SHA-1, SHA-224, SHA-256, SHA-384, SHA-512 y BLAKE2. Una función
hash transforma cualquier dato de entrada en una cadena de longitud fija
llamada ``digest'' o resumen, de manera que un cambio mínimo en la
entrada produce un resumen completamente diferente. Esta propiedad,
conocida como efecto avalancha, es fundamental para verificar la
integridad de archivos y contraseñas (Python Software Foundation, 2024).

\paragraph{4.1.3. Características
principales}\label{caracteruxedsticas-principales}

\begin{itemize}
\tightlist
\item
  Incluye algoritmos como MD5, SHA-1, SHA-256, SHA-512 y BLAKE2 sin
  instalación adicional.
\item
  Permite calcular hashes de cadenas de texto, archivos y flujos de
  datos de forma eficiente.
\item
  Es compatible con todas las plataformas donde corre Python (Windows,
  Linux, macOS).
\item
  Nivel de dificultad: básico. Su interfaz es simple e intuitiva para
  principiantes.
\item
  Se integra con otras librerías de ciberseguridad como hmac y secrets.
\end{itemize}

\paragraph{4.1.4. Instalación e
importación}\label{instalaciuxf3n-e-importaciuxf3n}

hashlib es parte de la biblioteca estándar de Python. No requiere
instalación con pip. Está disponible de forma inmediata en Google Colab,
Kaggle y cualquier entorno Python.

\begin{Shaded}
\begin{Highlighting}[]
\CommentTok{\# No se necesita instalación}
\ImportTok{import}\NormalTok{ hashlib}
\end{Highlighting}
\end{Shaded}

\paragraph{4.1.5. Funciones o métodos
principales}\label{funciones-o-muxe9todos-principales}

\begin{longtable}[]{@{}
  >{\raggedright\arraybackslash}p{(\linewidth - 4\tabcolsep) * \real{0.3333}}
  >{\raggedright\arraybackslash}p{(\linewidth - 4\tabcolsep) * \real{0.3333}}
  >{\raggedright\arraybackslash}p{(\linewidth - 4\tabcolsep) * \real{0.3333}}@{}}
\toprule\noalign{}
\begin{minipage}[b]{\linewidth}\raggedright
Función / Método
\end{minipage} & \begin{minipage}[b]{\linewidth}\raggedright
¿Para qué sirve?
\end{minipage} & \begin{minipage}[b]{\linewidth}\raggedright
Ejemplo breve
\end{minipage} \\
\midrule\noalign{}
\endhead
\bottomrule\noalign{}
\endlastfoot
\texttt{hashlib.sha256()} & Crea un objeto hash SHA-256. &
\texttt{h\ =\ hashlib.sha256()} \\
\texttt{hashlib.md5()} & Crea un objeto hash MD5. No recomendado para
seguridad. & \texttt{h\ =\ hashlib.md5()} \\
\texttt{h.update(data)} & Alimenta datos al hash incrementalmente. &
\texttt{h.update(b\textquotesingle{}mensaje\textquotesingle{})} \\
\texttt{h.hexdigest()} & Retorna el hash en hexadecimal. &
\texttt{print(h.hexdigest())} \\
\texttt{h.digest()} & Retorna el hash como bytes. &
\texttt{h.digest()} \\
\texttt{hashlib.algorithms\_available} & Conjunto de algoritmos hash
disponibles. & \texttt{hashlib.algorithms\_available} \\
\texttt{hashlib.new(\textquotesingle{}sha512\textquotesingle{},\ data)}
& Crea un objeto hash por nombre de algoritmo. &
\texttt{hashlib.new(\textquotesingle{}sha512\textquotesingle{},\ b\textquotesingle{}dato\textquotesingle{})} \\
\end{longtable}

\paragraph{4.1.6. Ejemplo básico de uso}\label{ejemplo-buxe1sico-de-uso}

El siguiente fragmento, extraído del módulo \texttt{comparar\_hashes.py}
del repositorio del proyecto, calcula simultáneamente los hashes MD5,
SHA-1, SHA-256 y SHA-512 de dos archivos APK ---uno original y otro
envenenado--- para detectar diferencias causadas por la inyección de un
payload malicioso:

\begin{Shaded}
\begin{Highlighting}[]
\ImportTok{import}\NormalTok{ hashlib, os, time}

\NormalTok{DIR }\OperatorTok{=} \StringTok{"/home/alim/exposicion final/APKs"}
\NormalTok{ARCHIVOS }\OperatorTok{=}\NormalTok{ [}\StringTok{"CairosOriginal.apk"}\NormalTok{, }\StringTok{"CairosVenom.apk"}\NormalTok{]}

\KeywordTok{def}\NormalTok{ calcular\_hashes(ruta):}
\NormalTok{    algoritmos }\OperatorTok{=}\NormalTok{ \{}
        \StringTok{"MD5"}\NormalTok{: hashlib.md5(),}
        \StringTok{"SHA{-}1"}\NormalTok{: hashlib.sha1(),}
        \StringTok{"SHA{-}256"}\NormalTok{: hashlib.sha256(),}
        \StringTok{"SHA{-}512"}\NormalTok{: hashlib.sha512(),}
\NormalTok{    \}}
    \ControlFlowTok{with} \BuiltInTok{open}\NormalTok{(ruta, }\StringTok{"rb"}\NormalTok{) }\ImportTok{as}\NormalTok{ f:}
        \ControlFlowTok{for}\NormalTok{ chunk }\KeywordTok{in} \BuiltInTok{iter}\NormalTok{(}\KeywordTok{lambda}\NormalTok{: f.read(}\DecValTok{4096}\NormalTok{), }\StringTok{b""}\NormalTok{):}
            \ControlFlowTok{for}\NormalTok{ h }\KeywordTok{in}\NormalTok{ algoritmos.values():}
\NormalTok{                h.update(chunk)}
    \ControlFlowTok{return}\NormalTok{ \{nombre: obj.hexdigest() }\ControlFlowTok{for}\NormalTok{ nombre, obj }\KeywordTok{in}\NormalTok{ algoritmos.items()\}}
\end{Highlighting}
\end{Shaded}

Este diseño es eficiente porque lee el archivo en bloques de 4096 bytes,
evitando cargar el archivo completo en memoria (los APK analizados
superan los 52 MB). Además, calcular los cuatro hashes en una sola
pasada reduce el tiempo de procesamiento.

\textbf{Demostración del efecto avalancha:} Una vez calculados los
hashes SHA-256 de ambos APK, el script aplica una operación XOR bit a
bit entre los dos digests para cuantificar el porcentaje de bits que
cambiaron:

\begin{Shaded}
\begin{Highlighting}[]
\NormalTok{h1 }\OperatorTok{=} \BuiltInTok{bytes}\NormalTok{.fromhex(original[}\StringTok{"SHA{-}256"}\NormalTok{])}
\NormalTok{h2 }\OperatorTok{=} \BuiltInTok{bytes}\NormalTok{.fromhex(envenenado[}\StringTok{"SHA{-}256"}\NormalTok{])}
\NormalTok{bits\_distintos }\OperatorTok{=} \BuiltInTok{sum}\NormalTok{(}\BuiltInTok{bin}\NormalTok{(a }\OperatorTok{\^{}}\NormalTok{ b).count(}\StringTok{\textquotesingle{}1\textquotesingle{}}\NormalTok{) }\ControlFlowTok{for}\NormalTok{ a, b }\KeywordTok{in} \BuiltInTok{zip}\NormalTok{(h1, h2))}
\NormalTok{total\_bits }\OperatorTok{=} \BuiltInTok{len}\NormalTok{(h1) }\OperatorTok{*} \DecValTok{8}
\NormalTok{porcentaje }\OperatorTok{=}\NormalTok{ (bits\_distintos }\OperatorTok{/}\NormalTok{ total\_bits) }\OperatorTok{*} \DecValTok{100}
\end{Highlighting}
\end{Shaded}

El resultado obtenido fue de 130 bits diferentes de 256 (50.8\%),
confirmando el efecto avalancha: la inyección del payload de Metasploit
alteró aproximadamente la mitad de los bits del hash, haciendo imposible
predecir el nuevo valor a partir del original.

\paragraph{4.1.7. Aplicación en contexto
real}\label{aplicaciuxf3n-en-contexto-real}

En el marco de esta investigación se aplicó hashlib al análisis forense
de dos aplicaciones Android reales: \texttt{CairosOriginal.apk}
(52,400,958 bytes) y \texttt{CairosVenom.apk} (52,431,637 bytes). El
segundo archivo fue generado mediante MsfVenom, inyectando un payload
meterpreter\_reverse\_tcp en el APK original. El script
\texttt{comparar\_hashes.py} calculó los cuatro hashes de cada archivo,
obteniendo valores completamente diferentes en todos los algoritmos
(MD5, SHA-1, SHA-256 y SHA-512), lo que demuestra que la inyección del
payload ---que solo agrega 30,679 bytes al archivo--- produjo un cambio
radical en los resúmenes criptográficos. Este mismo principio es
utilizado por plataformas como VirusTotal para detectar archivos
modificados maliciosamente.

\paragraph{4.1.8. Ventajas y
limitaciones}\label{ventajas-y-limitaciones}

\begin{longtable}[]{@{}
  >{\raggedright\arraybackslash}p{(\linewidth - 2\tabcolsep) * \real{0.5000}}
  >{\raggedright\arraybackslash}p{(\linewidth - 2\tabcolsep) * \real{0.5000}}@{}}
\toprule\noalign{}
\begin{minipage}[b]{\linewidth}\raggedright
Ventajas
\end{minipage} & \begin{minipage}[b]{\linewidth}\raggedright
Limitaciones
\end{minipage} \\
\midrule\noalign{}
\endhead
\bottomrule\noalign{}
\endlastfoot
Incluida en Python: no requiere instalación. & MD5 y SHA-1 son
considerados inseguros para uso criptográfico crítico. \\
Interfaz simple y bien documentada. & No cifra datos: solo genera un
resumen unidireccional (no es reversible). \\
Compatible con múltiples algoritmos de hash. & Para operaciones de
contraseñas, se recomienda usar bcrypt o argon2 en lugar de hashlib. \\
Eficiente para procesar archivos de gran tamaño mediante bloques. & No
gestiona claves ni certificados digitales por sí sola. \\
\end{longtable}

\subsubsection{4.2. LIBRERÍA N.º 2:
yara-python}\label{libreruxeda-n.uxba-2-yara-python}

\paragraph{4.2.1. Nombre de la
librería}\label{nombre-de-la-libreruxeda-1}

Nombre oficial: yara-python. Se instala con
\texttt{pip\ install\ yara-python} y se importa como:
\texttt{import\ yara}. Es la interfaz Python de YARA (Yet Another
Ridiculous Acronym), la herramienta de detección de malware más
utilizada en entornos de análisis forense y respuesta a incidentes.

\paragraph{4.2.2. Definición de la
librería}\label{definiciuxf3n-de-la-libreruxeda-1}

yara-python es el módulo de Python que expone la API de YARA, un motor
de detección de patrones diseñado para identificar y clasificar muestras
de malware a partir de reglas definidas por el analista. YARA permite
describir características de familias de malware mediante reglas que
combinan cadenas de texto, patrones hexadecimales y condiciones lógicas.
Fue creado originalmente por Víctor Manuel Álvarez y es mantenido
actualmente por VirusTotal (Google) (VirusTotal, 2024).

\paragraph{4.2.3. Características
principales}\label{caracteruxedsticas-principales-1}

\begin{itemize}
\tightlist
\item
  Permite escribir reglas de detección basadas en patrones de texto,
  bytes hexadecimales y expresiones regulares.
\item
  Analiza archivos, directorios y procesos en memoria en busca de
  coincidencias con las reglas definidas.
\item
  Nivel de dificultad: intermedio. Requiere aprender la sintaxis de
  reglas YARA.
\item
  Compatible con Python 3 y se integra con herramientas como Volatility
  y frameworks de análisis de malware.
\item
  Es la librería base de plataformas profesionales como VirusTotal,
  Cuckoo Sandbox y MISP.
\end{itemize}

\paragraph{4.2.4. Instalación e
importación}\label{instalaciuxf3n-e-importaciuxf3n-1}

\begin{Shaded}
\begin{Highlighting}[]
\CommentTok{\# Instalación}
\NormalTok{pip install yara}\OperatorTok{{-}}\NormalTok{python}

\CommentTok{\# Importación}
\ImportTok{import}\NormalTok{ yara}
\end{Highlighting}
\end{Shaded}

\paragraph{4.2.5. Funciones o métodos
principales}\label{funciones-o-muxe9todos-principales-1}

\begin{longtable}[]{@{}
  >{\raggedright\arraybackslash}p{(\linewidth - 4\tabcolsep) * \real{0.3333}}
  >{\raggedright\arraybackslash}p{(\linewidth - 4\tabcolsep) * \real{0.3333}}
  >{\raggedright\arraybackslash}p{(\linewidth - 4\tabcolsep) * \real{0.3333}}@{}}
\toprule\noalign{}
\begin{minipage}[b]{\linewidth}\raggedright
Función / Método
\end{minipage} & \begin{minipage}[b]{\linewidth}\raggedright
¿Para qué sirve?
\end{minipage} & \begin{minipage}[b]{\linewidth}\raggedright
Ejemplo breve
\end{minipage} \\
\midrule\noalign{}
\endhead
\bottomrule\noalign{}
\endlastfoot
\texttt{yara.compile(source=...)} & Compila una regla YARA desde una
cadena. & \texttt{r\ =\ yara.compile(source=s)} \\
\texttt{yara.compile(filepath=...)} & Compila reglas desde un archivo
.yar. &
\texttt{r\ =\ yara.compile(filepath=\textquotesingle{}r.yar\textquotesingle{})} \\
\texttt{regla.match(filepath=...)} & Escanea un archivo y retorna
coincidencias. &
\texttt{matches\ =\ r.match(filepath=\textquotesingle{}f.exe\textquotesingle{})} \\
\texttt{regla.match(data=...)} & Escanea datos en memoria (bytes). &
\texttt{r.match(data=contenido)} \\
\texttt{matches{[}i{]}.rule} & Nombre de la regla que coincidió. &
\texttt{print(matches{[}0{]}.rule)} \\
\texttt{matches{[}i{]}.strings} & Cadenas que activaron la regla. &
\texttt{print(matches{[}0{]}.strings)} \\
\end{longtable}

\paragraph{4.2.6. Ejemplo básico de
uso}\label{ejemplo-buxe1sico-de-uso-1}

El siguiente fragmento, tomado de \texttt{detector\_yara.py}, define
dos reglas YARA para detectar payloads de Metasploit dentro del archivo
classes.dex de un APK. Las reglas se compilan y posteriormente se aplican
al DEX extraído del paquete:

\begin{Shaded}
\begin{Highlighting}[]
\ImportTok{import}\NormalTok{ yara}

\NormalTok{REGLAS\_DEX }\OperatorTok{=}\NormalTok{ yara.}\BuiltInTok{compile}\NormalTok{(source}\OperatorTok{=}\StringTok{\textquotesingle{}\textquotesingle{}\textquotesingle{}}
\StringTok{    rule dex\_metasploit\_payload \{}
\StringTok{        meta:}
\StringTok{            description = "Detecta payload de Metasploit en DEX"}
\StringTok{            severity = "critical"}
\StringTok{        strings:}
\StringTok{            \$a = "Ljava/lang/reflect/Method;"}
\StringTok{            \$b = "Ljava/net/URLConnection;"}
\StringTok{            \$c = "addRequestProperty"}
\StringTok{            \$d = ";->startService(Landroid/content/Context;)V"}
\StringTok{            \$e = ";->openConnection()Ljava/net/URLConnection;"}
\StringTok{        condition:}
\StringTok{            uint32(0) == 0x0A786564 and}
\StringTok{            4 of them}
\StringTok{    \}}
\StringTok{\textquotesingle{}\textquotesingle{}\textquotesingle{}}\NormalTok{)}

\CommentTok{\# Extraer DEX del APK y escanear con YARA}
\ControlFlowTok{with}\NormalTok{ zipfile.ZipFile(ruta\_apk, }\StringTok{\textquotesingle{}r\textquotesingle{}}\NormalTok{) }\ImportTok{as}\NormalTok{ z:}
    \ControlFlowTok{for}\NormalTok{ name }\KeywordTok{in}\NormalTok{ z.namelist():}
        \ControlFlowTok{if}\NormalTok{ name.endswith(}\StringTok{\textquotesingle{}.dex\textquotesingle{}}\NormalTok{):}
\NormalTok{            dex\_data }\OperatorTok{=}\NormalTok{ z.read(name)}
            \ControlFlowTok{break}

\NormalTok{matches }\OperatorTok{=}\NormalTok{ REGLAS\_DEX.match(data}\OperatorTok{=}\NormalTok{dex\_data)}
\ControlFlowTok{if}\NormalTok{ matches:}
    \BuiltInTok{print}\NormalTok{(}\StringTok{\textquotesingle{}ALERTA: Payload detectado\textquotesingle{}}\NormalTok{)}
\end{Highlighting}
\end{Shaded}

La regla \texttt{dex\_metasploit\_payload} busca cinco indicadores
característicos de un RAT Android: uso de reflexión
(\texttt{java/lang/reflect/Method}), conexiones de red
(\texttt{java/net/URLConnection}), envío de cabeceras HTTP
(\texttt{addRequestProperty}), inicio de servicios remotos y apertura de
conexiones URL. La condición exige que al menos cuatro de estos cinco
patrones estén presentes, y además verifica que el archivo comience con
el magic number de un DEX válido (\texttt{0x0A786564}).

\textbf{Análisis combinado con Androguard:} Además del escaneo con
YARA, el detector implementa heurísticas complementarias: cuenta los
permisos peligrosos declarados en el AndroidManifest.xml (24 en el APK
envenenado frente a 2 en el original), identifica subpaquetes inyectados
de nombre aleatorio corto (como \texttt{com.example.cairos.dfwto}), y
compara el tamaño del DEX (993 KB vs 812 KB). Estos indicadores se
combinan en una puntuación de riesgo que determina el veredicto final.

\paragraph{4.2.7. Aplicación en contexto
real}\label{aplicaciuxf3n-en-contexto-real-1}

El detector \texttt{detector\_yara.py} desarrollado en este proyecto
analizó tres APK del repositorio: \texttt{CairosOriginal.apk},
\texttt{CairosVenom.apk} y \texttt{payload.apk}. El sistema identificó
correctamente los dos APK envenenados mediante la combinación de reglas
YARA sobre el DEX y heurísticas sobre el manifiesto Android. El APK
original fue clasificado como LIMPIO (2 permisos, sin subpaquetes),
mientras que el APK envenenado fue clasificado como ENVENENADO (24
permisos, 14 de ellos peligrosos, un subpaquete inyectado con nombre
aleatorio de 5 caracteres). Este enfoque multicapa ---YARA para firmas
binarias y Androguard para análisis estructural--- es el mismo que
emplean las soluciones antimalware comerciales para Android.

\paragraph{4.2.8. Ventajas y
limitaciones}\label{ventajas-y-limitaciones-1}

\begin{longtable}[]{@{}
  >{\raggedright\arraybackslash}p{(\linewidth - 2\tabcolsep) * \real{0.5000}}
  >{\raggedright\arraybackslash}p{(\linewidth - 2\tabcolsep) * \real{0.5000}}@{}}
\toprule\noalign{}
\begin{minipage}[b]{\linewidth}\raggedright
Ventajas
\end{minipage} & \begin{minipage}[b]{\linewidth}\raggedright
Limitaciones
\end{minipage} \\
\midrule\noalign{}
\endhead
\bottomrule\noalign{}
\endlastfoot
Es el estándar industrial para detección de malware basada en reglas. &
Requiere aprender la sintaxis YARA para escribir reglas efectivas. \\
Altamente expresiva: combina strings, bytes, regex y lógica booleana. &
Las reglas mal escritas pueden generar falsos positivos o negativos. \\
Compatible con miles de reglas públicas disponibles en repositorios como
GitHub. & No detecta malware polimórfico avanzado que modifica sus
patrones constantemente. \\
Permite escanear archivos, procesos en memoria y flujos de datos. &
Requiere actualización constante de las reglas ante nuevas amenazas. \\
\end{longtable}

\subsubsection{4.3. LIBRERÍA N.º 3:
impacket}\label{libreruxeda-n.uxba-3-impacket}

\paragraph{4.3.1. Nombre de la
librería}\label{nombre-de-la-libreruxeda-2}

Nombre oficial: impacket. Se instala con \texttt{pip\ install\ impacket}
y se accede a sus módulos mediante importaciones específicas según el
protocolo: \texttt{from\ impacket.smbconnection\ import\ SMBConnection},
\texttt{from\ impacket\ import\ smb,\ ntlm}, entre otros.

\paragraph{4.3.2. Definición de la
librería}\label{definiciuxf3n-de-la-libreruxeda-2}

impacket es una colección de clases Python para trabajar con protocolos
de red de bajo nivel, especialmente los utilizados en entornos Windows:
SMB, MSRPC, NTLM, Kerberos, LDAP, MSSQL, WMI y DCOM. Fue desarrollada
originalmente por SecureAuth Corporation y actualmente es mantenida por
la comunidad a través de Fortra. Es ampliamente utilizada en auditorías
de seguridad, análisis forense de red y laboratorios de seguridad
ofensiva y defensiva (Fortra, 2024).

\textbf{IMPORTANTE:} impacket debe utilizarse exclusivamente en entornos
propios, laboratorios controlados o con autorización escrita del
propietario del sistema. Su uso no autorizado en sistemas de terceros
constituye un delito informático.

\paragraph{4.3.3. Características
principales}\label{caracteruxedsticas-principales-2}

\begin{itemize}
\tightlist
\item
  Implementa protocolos como SMB1/2/3, Kerberos, NTLM, LDAP, MSSQL, WMI
  y DCOM en Python puro.
\item
  Permite construir, enviar y analizar paquetes de red a nivel de
  protocolo.
\item
  Nivel de dificultad: avanzado. Requiere conocimiento de protocolos de
  red Windows.
\item
  Es la base de herramientas de seguridad reconocidas como BloodHound,
  crackmapexec y secretsdump.
\item
  Incluye más de 20 scripts de ejemplo listos para usar en auditorías de
  seguridad.
\end{itemize}

\paragraph{4.3.4. Instalación e
importación}\label{instalaciuxf3n-e-importaciuxf3n-2}

\begin{Shaded}
\begin{Highlighting}[]
\CommentTok{\# Instalación}
\NormalTok{pip install impacket}

\CommentTok{\# Importaciones según el módulo necesario}
\ImportTok{from}\NormalTok{ impacket.smbconnection }\ImportTok{import}\NormalTok{ SMBConnection}
\ImportTok{from}\NormalTok{ impacket }\ImportTok{import}\NormalTok{ smb, ntlm}
\ImportTok{from}\NormalTok{ impacket.krb5 }\ImportTok{import}\NormalTok{ kerberosv5}
\end{Highlighting}
\end{Shaded}

\paragraph{4.3.5. Funciones o métodos
principales}\label{funciones-o-muxe9todos-principales-2}

\begin{longtable}[]{@{}
  >{\raggedright\arraybackslash}p{(\linewidth - 4\tabcolsep) * \real{0.3333}}
  >{\raggedright\arraybackslash}p{(\linewidth - 4\tabcolsep) * \real{0.3333}}
  >{\raggedright\arraybackslash}p{(\linewidth - 4\tabcolsep) * \real{0.3333}}@{}}
\toprule\noalign{}
\begin{minipage}[b]{\linewidth}\raggedright
Módulo / Función
\end{minipage} & \begin{minipage}[b]{\linewidth}\raggedright
¿Para qué sirve?
\end{minipage} & \begin{minipage}[b]{\linewidth}\raggedright
Ejemplo breve
\end{minipage} \\
\midrule\noalign{}
\endhead
\bottomrule\noalign{}
\endlastfoot
\texttt{SMBConnection()} & Establece conexión SMB a un servidor. &
\texttt{conn\ =\ SMBConnection(ip,\ ip)} \\
\texttt{conn.login(user,\ pass)} & Autentica la conexión con
credenciales. &
\texttt{conn.login(\textquotesingle{}user\textquotesingle{},\ \textquotesingle{}pass\textquotesingle{})} \\
\texttt{conn.listShares()} & Lista recursos compartidos del servidor. &
\texttt{shares\ =\ conn.listShares()} \\
\texttt{ntlm.getNTLMSSPType1()} & Construye mensaje NTLM Tipo 1. &
\texttt{msg\ =\ getNTLMSSPType1()} \\
\texttt{smb.SMB.get\_server\_name()} & Obtiene el nombre del servidor
SMB. & \texttt{server\ =\ get\_server\_name()} \\
\texttt{kerberosv5.getKerberosTGT()} & Solicita un TGT de Kerberos. &
\texttt{tgt\ =\ getKerberosTGT(...)} \\
\end{longtable}

\paragraph{4.3.6. Limitaciones para la documentación
práctica}\label{limitaciones-para-la-documentaciuxf3n-pruxe1ctica}

Durante el desarrollo de esta monografía se intentó generar una práctica
funcional completa con impacket, similar a la realizada con hashlib y
yara-python. Sin embargo, se encontraron las siguientes barreras
técnicas y legales que imposibilitaron la ejecución de pruebas reales:

\begin{enumerate}
\def\labelenumi{\arabic{enumi}.}
\tightlist
\item
  \textbf{Dependencia de un dominio Active Directory real:} Los módulos
  Kerberos de impacket requieren un Key Distribution Center (KDC)
  funcional para solicitar y validar Ticket Granting Tickets (TGT). Sin
  un controlador de dominio Windows Server configurado con usuarios,
  grupos y políticas de autenticación, las funciones
  \texttt{getKerberosTGT()} y \texttt{getKerberosTGS()} retornan errores
  de conexión irrecuperables. Montar un laboratorio completo con Windows
  Server, roles de dominio y cuentas de servicio excede el alcance de
  este trabajo monográfico.
\item
  \textbf{Necesidad de credenciales privilegiadas:} Incluso en un entorno
  de laboratorio, los scripts de impacket como \texttt{secretsdump.py},
  \texttt{GetNPUsers.py} y \texttt{wmiexec.py} requieren autenticación
  con cuentas de dominio que posean privilegios específicos (e.g.,
  permisos de replicación de directorio para DCSync). La generación de
  estas credenciales implica la administración completa de un dominio,
  actividad que requiere autorización institucional formal.
\item
  \textbf{Complejidad del código base:} impacket implementa más de 20
  módulos que abarcan protocolos completos a nivel de paquete. Cada
  módulo encapsula estructuras ASN.1, mecanismos de negociación GSS-API,
  cifrado RC4-HMAC y AES para Kerberos, y parsing de PAC (Privilege
  Attribute Certificate). Documentar un flujo completo requeriría
  desensamblar y explicar miles de líneas de código que operan a niveles
  del stack OSI que exceden los objetivos de esta investigación.
\item
  \textbf{Implicaciones éticas y legales:} La ejecución de herramientas
  de impacket contra sistemas que no son de propiedad del equipo
  investigador constituiría un delito informático tipificado en la
  legislación peruana (Ley N° 30096, Ley de Delitos Informáticos). Por
  ello, la documentación se limita al análisis conceptual de la
  arquitectura de la librería y sus métodos principales, sin ejecución
  contra infraestructura real.
\end{enumerate}

No obstante, se reconoce que impacket es una herramienta de valor
incalculable en auditorías de seguridad profesionales. Su inclusión en
esta monografía responde al objetivo de documentar su existencia,
arquitectura y aplicaciones teóricas, dejando la práctica operativa para
entornos controlados con las autorizaciones correspondientes.

\paragraph{4.3.7. Aplicación en contexto
real}\label{aplicaciuxf3n-en-contexto-real-2}

Un auditor de seguridad contratado por una empresa utiliza impacket en
un entorno controlado para verificar si los servidores Windows de la
organización tienen recursos SMB expuestos sin autenticación adecuada, o
si el protocolo NTLM está configurado de forma vulnerable. Estas
pruebas, realizadas con autorización formal, permiten identificar
vectores de ataque antes de que actores externos los exploten,
contribuyendo directamente a la postura de seguridad de la organización.

\paragraph{4.3.8. Ventajas y
limitaciones}\label{ventajas-y-limitaciones-2}

\begin{longtable}[]{@{}
  >{\raggedright\arraybackslash}p{(\linewidth - 2\tabcolsep) * \real{0.5000}}
  >{\raggedright\arraybackslash}p{(\linewidth - 2\tabcolsep) * \real{0.5000}}@{}}
\toprule\noalign{}
\begin{minipage}[b]{\linewidth}\raggedright
Ventajas
\end{minipage} & \begin{minipage}[b]{\linewidth}\raggedright
Limitaciones
\end{minipage} \\
\midrule\noalign{}
\endhead
\bottomrule\noalign{}
\endlastfoot
Implementación Python pura de protocolos de red Windows sin dependencias
externas. & Nivel avanzado: requiere sólidos conocimientos de protocolos
de red. \\
Base de herramientas profesionales ampliamente utilizadas en la
industria. & Su uso fuera de entornos autorizados puede constituir un
delito. \\
Permite analizar y entender protocolos como SMB y Kerberos a nivel de
paquete. & La documentación oficial es técnica y puede ser difícil para
principiantes. \\
Actualización constante para cubrir nuevas versiones de protocolos. &
Requiere un entorno de laboratorio (máquina virtual) para practicar de
forma segura. \\
\end{longtable}

\subsection{CAPÍTULO V. CUADRO COMPARATIVO DE LAS
LIBRERÍAS}\label{capuxedtulo-v.-cuadro-comparativo-de-las-libreruxedas}

\begin{longtable}[]{@{}
  >{\raggedright\arraybackslash}p{(\linewidth - 6\tabcolsep) * \real{0.2500}}
  >{\raggedright\arraybackslash}p{(\linewidth - 6\tabcolsep) * \real{0.2500}}
  >{\raggedright\arraybackslash}p{(\linewidth - 6\tabcolsep) * \real{0.2500}}
  >{\raggedright\arraybackslash}p{(\linewidth - 6\tabcolsep) * \real{0.2500}}@{}}
\toprule\noalign{}
\begin{minipage}[b]{\linewidth}\raggedright
Criterio
\end{minipage} & \begin{minipage}[b]{\linewidth}\raggedright
hashlib
\end{minipage} & \begin{minipage}[b]{\linewidth}\raggedright
yara-python
\end{minipage} & \begin{minipage}[b]{\linewidth}\raggedright
impacket
\end{minipage} \\
\midrule\noalign{}
\endhead
\bottomrule\noalign{}
\endlastfoot
Área de aplicación & Verificación de integridad criptográfica. &
Detección y clasificación de malware. & Análisis y simulación de
protocolos de red. \\
Uso principal & Calcular hashes MD5, SHA-256, SHA-512 de archivos y
datos. & Escanear archivos y memoria con reglas de patrones YARA. &
Implementar y analizar protocolos SMB, NTLM, Kerberos, LDAP. \\
Nivel de dificultad & Básico & Intermedio & Avanzado \\
Tipo de datos que maneja & Strings, bytes, archivos de cualquier tipo. &
Archivos, procesos en memoria, flujos de bytes. & Paquetes de red,
sesiones SMB, tickets Kerberos. \\
Funciones más importantes & \texttt{sha256()}, \texttt{md5()},
\texttt{update()}, \texttt{hexdigest()} & \texttt{compile()},
\texttt{match()}, \texttt{strings}, \texttt{meta} &
\texttt{SMBConnection}, \texttt{login()}, \texttt{listShares()}, NTLM \\
Ventajas principales & Estándar, sin instalación, interfaz simple. &
Estándar industrial en análisis de malware. & Implementa protocolos
reales en Python puro. \\
Limitaciones & No cifra, solo genera resúmenes; MD5/SHA-1 inseguros. &
Requiere aprender sintaxis YARA; no detecta malware polimórfico
avanzado. & Nivel avanzado; uso solo en entornos autorizados. \\
Ejemplo de uso real & Verificar integridad de instaladores de software.
& Escanear adjuntos de correo en busca de malware. & Auditar
configuraciones SMB/Kerberos en red interna. \\
Relación con la práctica & Cálculo y comparación de hashes de APKs de 52 MB; efecto avalancha con XOR bit a bit. &
Detección de payload Metasploit en DEX + análisis de permisos y subpaquetes con Androguard. &
Documentación teórica; no se ejecutó práctica por requerir dominio Active Directory y credenciales autorizadas. \\
\end{longtable}

\subsubsection{Análisis interpretativo del cuadro
comparativo}\label{anuxe1lisis-interpretativo-del-cuadro-comparativo}

Al comparar las tres librerías, se observa una progresión natural en
términos de complejidad y alcance dentro de la ciberseguridad defensiva.
hashlib resultó la más accesible para principiantes por su interfaz
directa y su inclusión en la biblioteca estándar, siendo ideal como
punto de entrada al estudio de la criptografía aplicada. En la práctica
realizada, procesó archivos de 52 MB en 0.37 segundos calculando cuatro
hashes simultáneamente, y la comparación XOR bit a bit del SHA-256
cuantificó el efecto avalancha en 50.8\%.

yara-python exige un nivel intermedio de conocimiento, especialmente en
la sintaxis de reglas YARA, pero ofrece una capacidad de detección de
amenazas que se acerca a los entornos profesionales reales. En este
proyecto, se demostró que la combinación de reglas YARA sobre el DEX con
heurísticas de Androguard (conteo de permisos, detección de subpaquetes
inyectados) produce una clasificación precisa de APKs envenenados.

impacket es la herramienta de mayor complejidad técnica, orientada a
profesionales con conocimientos consolidados de redes, protocolos
Windows y administración de dominios Active Directory. Su documentación
práctica no pudo ser completada porque los módulos Kerberos requieren un
KDC funcional, los scripts de auditoría necesitan credenciales
privilegiadas de dominio, y la ejecución sin autorización constituiría
un delito informático. No obstante, se documentó su arquitectura,
métodos principales y aplicaciones teóricas.

Las tres librerías son complementarias: hashlib garantiza que los
archivos no han sido modificados, yara-python detecta si contienen
código malicioso, e impacket permite analizar cómo se comunican los
sistemas en red. Juntas, conforman una cadena de análisis defensivo
coherente, aunque impacket requiere condiciones de laboratorio que
exceden el alcance típico de un trabajo monográfico de pregrado.

\subsection{CAPÍTULO VI. PARTE PRÁCTICA: ANÁLISIS FORENSE DE
APKs}\label{capuxedtulo-vi.-parte-pruxe1ctica-en-google-colab}

\subsubsection{6.1. Plataforma utilizada}\label{plataforma-utilizada}

Se utilizó un entorno local Linux con Python 3.14, un entorno virtual
(\texttt{venv}) con las dependencias yara-python y androguard, y las
herramientas del sistema apktool, msfvenom y Metasploit Framework para
la generación del APK envenenado. Para la versión de Google Colab, se
encuentran disponibles dos notebooks: uno para hashlib y otro para
yara-python (ver Anexo 4). Los APK utilizados y todo el código fuente
están publicados en el repositorio GitHub del proyecto.

\subsubsection{6.2. Descripción del problema
práctico}\label{descripciuxf3n-del-problema-pruxe1ctico}

El problema práctico consiste en determinar, mediante herramientas
Python, si un archivo APK ha sido manipulado con un payload malicioso
inyectado por MsfVenom (Metasploit). Se trabaja con tres archivos:
\texttt{CairosOriginal.apk} (aplicación Flutter legítima, 52.4 MB),
\texttt{CairosVenom.apk} (el mismo APK con un payload
meterpreter\_reverse\_tcp inyectado, 52.4 MB) y \texttt{payload.apk}
(APK standalone generado directamente por msfvenom, 10 KB). El análisis
se realiza en dos fases: (1) comparación de hashes criptográficos para
verificar la integridad, y (2) detección automatizada mediante reglas
YARA y heurísticas estructurales con Androguard.

\subsubsection{6.3. Datos utilizados}\label{datos-utilizados}

\begin{itemize}
\tightlist
\item
  \texttt{CairosOriginal.apk} (52,400,958 bytes): APK legítimo original.
\item
  \texttt{CairosVenom.apk} (52,431,637 bytes): APK original con payload
  Metasploit inyectado.
\item
  \texttt{payload.apk} (10,246 bytes): APK standalone de msfvenom.
\end{itemize}

Todos los APK se encuentran en la carpeta \texttt{APKs/} del
repositorio.

\subsubsection{6.4. Código desarrollado}\label{cuxf3digo-desarrollado}

\textbf{Fase 1: Comparación de hashes con hashlib}

El script \texttt{comparar\_hashes.py} calcula MD5, SHA-1, SHA-256 y
SHA-512 de los dos APK principales y compara bit a bit los resúmenes:

\begin{Shaded}
\begin{Highlighting}[]
\ImportTok{import}\NormalTok{ hashlib, os, time}

\NormalTok{DIR }\OperatorTok{=} \StringTok{"/home/alim/exposicion final/APKs"}
\NormalTok{ARCHIVOS }\OperatorTok{=}\NormalTok{ [}\StringTok{"CairosOriginal.apk"}\NormalTok{, }\StringTok{"CairosVenom.apk"}\NormalTok{]}

\KeywordTok{def}\NormalTok{ calcular\_hashes(ruta):}
\NormalTok{    algoritmos }\OperatorTok{=}\NormalTok{ \{}
        \StringTok{"MD5"}\NormalTok{: hashlib.md5(),}
        \StringTok{"SHA{-}1"}\NormalTok{: hashlib.sha1(),}
        \StringTok{"SHA{-}256"}\NormalTok{: hashlib.sha256(),}
        \StringTok{"SHA{-}512"}\NormalTok{: hashlib.sha512(),}
\NormalTok{    \}}
    \ControlFlowTok{with} \BuiltInTok{open}\NormalTok{(ruta, }\StringTok{"rb"}\NormalTok{) }\ImportTok{as}\NormalTok{ f:}
        \ControlFlowTok{for}\NormalTok{ chunk }\KeywordTok{in} \BuiltInTok{iter}\NormalTok{(}\KeywordTok{lambda}\NormalTok{: f.read(}\DecValTok{4096}\NormalTok{), }\StringTok{b""}\NormalTok{):}
            \ControlFlowTok{for}\NormalTok{ h }\KeywordTok{in}\NormalTok{ algoritmos.values():}
\NormalTok{                h.update(chunk)}
    \ControlFlowTok{return}\NormalTok{ \{nombre: obj.hexdigest() }\ControlFlowTok{for}\NormalTok{ nombre, obj }\KeywordTok{in}\NormalTok{ algoritmos.items()\}}
\end{Highlighting}
\end{Shaded}

\textbf{Fase 2: Detección con YARA + Androguard}

El script \texttt{detector\_yara.py} implementa un sistema de puntuación
que combina reglas YARA sobre el DEX con heurísticas del manifiesto
Android:

\begin{Shaded}
\begin{Highlighting}[]
\ImportTok{from}\NormalTok{ androguard.core.apk }\ImportTok{import}\NormalTok{ APK}

\NormalTok{PERMISOS\_PELIGROSOS }\OperatorTok{=}\NormalTok{ [}
    \StringTok{"android.permission.READ\_SMS"}\NormalTok{,}
    \StringTok{"android.permission.SEND\_SMS"}\NormalTok{,}
    \StringTok{"android.permission.CAMERA"}\NormalTok{,}
    \StringTok{"android.permission.RECORD\_AUDIO"}\NormalTok{,}
    \StringTok{"android.permission.READ\_CALL\_LOG"}\NormalTok{,}
\NormalTok{]}

\NormalTok{apk }\OperatorTok{=}\NormalTok{ APK(ruta)}
\NormalTok{permisos }\OperatorTok{=}\NormalTok{ apk.get\_permissions()}
\NormalTok{peligrosos }\OperatorTok{=}\NormalTok{ [p }\ControlFlowTok{for}\NormalTok{ p }\KeywordTok{in}\NormalTok{ permisos }\ControlFlowTok{if}\NormalTok{ p }\KeywordTok{in}\NormalTok{ PERMISOS\_PELIGROSOS]}

\CommentTok{\# Detectar subpaquete inyectado (nombre aleatorio corto)}
\NormalTok{main\_package }\OperatorTok{=}\NormalTok{ apk.get\_package()}
\NormalTok{subpaquetes }\OperatorTok{=}\NormalTok{ [p }\ControlFlowTok{for}\NormalTok{ p }\KeywordTok{in}\NormalTok{ paquetes}
              \ControlFlowTok{if}\NormalTok{ p.startswith(main\_package) }\KeywordTok{and}\NormalTok{ p }\OperatorTok{!=}\NormalTok{ main\_package]}
\end{Highlighting}
\end{Shaded}

\subsubsection{6.5. Explicación del
código}\label{explicaciuxf3n-del-cuxf3digo}

La Fase 1 utiliza lectura por bloques de 4 KB para procesar archivos de
más de 50 MB sin saturar la memoria RAM. Calcula cuatro algoritmos en
una sola pasada, optimizando el tiempo de ejecución (0.37 segundos para
52 MB). La comparación XOR bit a bit del SHA-256 demuestra
cuantitativamente el efecto avalancha.

La Fase 2 emplea un enfoque de detección por capas: (a) reglas YARA que
buscan firmas de código malicioso en el DEX (reflexión, conexiones de
red, inicio de servicios); (b) análisis del AndroidManifest.xml mediante
Androguard para contar permisos peligrosos declarados; (c) detección de
subpaquetes inyectados con nombres aleatorios de 5 caracteres
generados por MsfVenom; (d) comparación del tamaño del DEX (el payload
agrega 181,696 bytes). Cada indicador suma puntos a una puntuación de
riesgo; si el puntaje supera 5, el APK se clasifica como ENVENENADO.

\subsubsection{6.6. Resultados obtenidos}\label{resultados-obtenidos}

\begin{longtable}[]{@{}
  >{\raggedright\arraybackslash}p{(\linewidth - 6\tabcolsep) * \real{0.2500}}
  >{\raggedright\arraybackslash}p{(\linewidth - 6\tabcolsep) * \real{0.2500}}
  >{\raggedright\arraybackslash}p{(\linewidth - 6\tabcolsep) * \real{0.2500}}
  >{\raggedright\arraybackslash}p{(\linewidth - 6\tabcolsep) * \real{0.2500}}@{}}
\toprule\noalign{}
\begin{minipage}[b]{\linewidth}\raggedright
Indicador
\end{minipage} & \begin{minipage}[b]{\linewidth}\raggedright
CairosOriginal
\end{minipage} & \begin{minipage}[b]{\linewidth}\raggedright
CairosVenom
\end{minipage} & \begin{minipage}[b]{\linewidth}\raggedright
payload.apk
\end{minipage} \\
\midrule\noalign{}
\endhead
\bottomrule\noalign{}
\endlastfoot
Tamaño & 52,400,958 B & 52,431,637 B & 10,246 B \\
Permisos totales & 2 & 24 & 23 \\
Permisos peligrosos & 0 & 14 & 14 \\
Subpaquetes inyectados & 0 & 1 (dfwto) & 0 \\
Tamaño DEX & 812 KB & 993 KB & 20 KB \\
SHA-256 (primeros 16) & 6752e6831ed260a7 & fe1452b4900c9214 & --- \\
Veredicto & LIMPIO & ENVENENADO & ENVENENADO \\
Bits SHA-256 distintos & --- & 130/256 (50.8\%) & --- \\
\end{longtable}

\subsubsection{6.7. Evidencias de
ejecución}\label{evidencias-de-ejecuciuxf3n}

\begin{itemize}
\tightlist
\item
  Capturas de pantalla de la terminal con la ejecución de
  \texttt{comparar\_hashes.py} mostrando los cuatro hashes de cada APK.
\item
  Salida de \texttt{detector\_yara.py} clasificando los tres APK con su
  puntuación de riesgo e indicadores.
\item
  Notebook de Google Colab con la demostración de hashlib:\\
  \url{https://colab.research.google.com/drive/1oXyDNrLdfR4-8oltR1WURTp8IMK65TTF}
\item
  Notebook de Google Colab con la demostración de yara-python:\\
  \url{https://colab.research.google.com/drive/173T7NiHw52sU_ySdK0RQoJzm5KdplAmp}
\item
  Repositorio GitHub completo:\\
  \url{https://github.com/AlvaroHuacacSoto/hashlib-Jara}
\end{itemize}
\subsection{CAPÍTULO VII. CASO PRÁCTICO: DETECCIÓN DE APK
ENVENENADO}\label{capuxedtulo-vii.-caso-pruxe1ctico-en-contexto-real}

\subsubsection{7.1. Descripción del caso}\label{descripciuxf3n-del-caso}

Un desarrollador independiente descubre que una versión modificada de su
aplicación Android \texttt{Cairos.apk} está circulando en foros de
terceros. La versión legítima es una app Flutter de 52 MB que solo
requiere permisos de INTERNET y un permiso interno de broadcast. La
versión sospechosa tiene el mismo tamaño aproximado, pero al instalarse
solicita acceso a SMS, cámara, micrófono, contactos y registros de
llamadas. El desarrollador necesita determinar técnicamente si el APK
fue manipulado y cómo.

\subsubsection{7.2. Problema identificado}\label{problema-identificado}

\begin{itemize}
\tightlist
\item
  El APK circulante solicita 24 permisos (14 peligrosos), frente a los 2
  permisos de la versión legítima.
\item
  El hash criptográfico del archivo no coincide con el original
  publicado por el desarrollador.
\item
  Se sospecha la inyección de un RAT (Remote Access Trojan) mediante
  herramientas de código abierto como Metasploit/MsfVenom.
\item
  Es necesario un método automatizado y reproducible para detectar la
  manipulación sin depender de inspección manual.
\end{itemize}

\subsubsection{7.3. Solución propuesta con
Python}\label{soluciuxf3n-propuesta-con-python}

Se desarrollaron dos scripts Python que abordan el problema desde dos
ángulos complementarios:

\begin{itemize}
\tightlist
\item
  \textbf{comparar\_hashes.py:} calcula y compara los hashes MD5, SHA-1,
  SHA-256 y SHA-512 de ambos APK. La diferencia en los cuatro algoritmos
  confirma que el archivo fue alterado. El análisis XOR del SHA-256
  revela que el 50.8\% de los bits cambiaron, lo cual es consistente con
  la inyección de un payload de aproximadamente 30 KB dentro de un
  archivo de 52 MB.
\item
  \textbf{detector\_yara.py:} implementa un motor de detección que
  combina (a) reglas YARA compiladas para buscar patrones de RAT en el
  DEX (reflexión, conexiones URL, inicio de servicios remotos), (b)
  análisis del manifiesto mediante Androguard para contar permisos
  peligrosos y detectar subpaquetes inyectados con nombres aleatorios, y
  (c) comparación del tamaño del DEX. El sistema asigna una puntuación
  de riesgo y emite un veredicto (LIMPIO / SOSPECHOSO / ENVENENADO).
\end{itemize}

\subsubsection{7.4. Beneficios de la
solución}\label{beneficios-de-la-soluciuxf3n}

\begin{itemize}
\tightlist
\item
  El desarrollador puede verificar la integridad de su APK en segundos
  comparando el hash SHA-256 con el valor de referencia.
\item
  El detector YARA + Androguard identifica automáticamente APKs
  envenenados sin necesidad de ingeniería inversa manual.
\item
  El enfoque es reproducible: cualquier analista puede clonar el
  repositorio, instalar las dependencias y ejecutar los scripts sobre
  nuevos APKs sospechosos.
\item
  La metodología es aplicable a la detección de malware Android en
  tiendas de aplicaciones alternativas, repositorios de APK y campañas
  de phishing móvil.
\end{itemize}

\begin{longtable}[]{@{}
  >{\raggedright\arraybackslash}p{(\linewidth - 2\tabcolsep) * \real{0.5000}}
  >{\raggedright\arraybackslash}p{(\linewidth - 2\tabcolsep) * \real{0.5000}}@{}}
\toprule\noalign{}
\begin{minipage}[b]{\linewidth}\raggedright
Elemento del caso
\end{minipage} & \begin{minipage}[b]{\linewidth}\raggedright
Descripción
\end{minipage} \\
\midrule\noalign{}
\endhead
\bottomrule\noalign{}
\endlastfoot
Contexto & Desarrollador de app Flutter descubre versión alterada de su
APK en distribución no autorizada. \\
Actor principal & Desarrollador / Analista de seguridad móvil. \\
Problema & APK alterado con 24 permisos (14 peligrosos) vs 2 permisos
originales; hashes no coinciden. \\
Datos o archivos & \texttt{CairosOriginal.apk} (52.4 MB),
\texttt{CairosVenom.apk} (52.4 MB), \texttt{payload.apk} (10 KB). \\
Librerías & hashlib (verificación de integridad), yara-python + androguard
(detección automatizada), impacket (documentación teórica). \\
Herramienta de ataque identificada & MsfVenom (Metasploit Framework)
con payload android/meterpreter/reverse\_tcp. \\
Resultado obtenido & Detección exitosa: original clasificado LIMPIO,
envenenado clasificado ENVENENADO con 14 indicadores de peligro. \\
Repositorio & \url{https://github.com/AlvaroHuacacSoto/hashlib-Jara} \\
\end{longtable}

\subsection{CONCLUSIONES}\label{conclusiones}

\begin{enumerate}
\def\labelenumi{\arabic{enumi}.}
\tightlist
\item
  Las librerías hashlib, yara-python e impacket demuestran que Python es
  una plataforma robusta para implementar soluciones de ciberseguridad
  defensiva. Cada librería aborda una dimensión diferente del problema
  de seguridad: integridad (hashlib), detección (yara-python) y análisis
  de protocolos (impacket). El proyecto práctico desarrollado ---análisis
  forense de APKs envenenados con payload de Metasploit--- demostró la
  efectividad de las dos primeras en un escenario real, mientras que
  impacket fue documentada teóricamente por las barreras de autenticación
  que requiere.
\item
  El análisis de hashes con \texttt{comparar\_hashes.py} confirmó el
  efecto avalancha: los cuatro algoritmos (MD5, SHA-1, SHA-256, SHA-512)
  produjeron valores completamente diferentes entre el APK original y el
  envenenado, con 130 de 256 bits (50.8\%) distintos en SHA-256. Esto
  valida el uso de hashes criptográficos como primera línea de defensa
  contra manipulaciones de software.
\item
  El detector \texttt{detector\_yara.py} clasificó correctamente los
  tres APK: \texttt{CairosOriginal.apk} como LIMPIO (2 permisos),
  \texttt{CairosVenom.apk} como ENVENENADO (24 permisos, subpaquete
  inyectado \texttt{dfwto}) y \texttt{payload.apk} como ENVENENADO. La
  combinación de reglas YARA con heurísticas de Androguard demostró ser
  más efectiva que cualquiera de las técnicas por separado.
\item
  El caso práctico evidenció que herramientas de código abierto como
  MsfVenom pueden ser utilizadas para generar APKs maliciosos
  indetectables a simple vista (mismo ícono, nombre y tamaño similar),
  pero las librerías Python de ciberseguridad permiten identificarlos
  mediante análisis automatizado de su estructura interna.
\end{enumerate}

\subsection{RECOMENDACIONES}\label{recomendaciones}

\begin{itemize}
\tightlist
\item
  Clonar el repositorio
  \url{https://github.com/AlvaroHuacacSoto/hashlib-Jara} que contiene
  todos los scripts, APKs de prueba y este documento para replicar los
  experimentos.
\item
  Revisar los notebooks de Google Colab publicados como complemento
  interactivo a esta monografía.
\item
  Al usar hashlib para verificación de integridad, preferir SHA-256 o
  SHA-512 y evitar MD5/SHA-1 por sus vulnerabilidades de colisión
  conocidas.
\item
  Combinar reglas YARA con análisis estructural del manifiesto Android
  para reducir falsos negativos en la detección de malware móvil.
\item
  Para trabajar con impacket, montar un laboratorio controlado con
  Windows Server y Active Directory, y obtener autorización formal antes
  de ejecutar cualquier script contra infraestructura que no sea propia.
\item
  Mantener actualizadas las reglas YARA incorporando firmas de
  repositorios comunitarios como YARA-Rules de GitHub y las colecciones
  de VirusTotal.
\end{itemize}

\subsection{REFERENCIAS
BIBLIOGRÁFICAS}\label{referencias-bibliogruxe1ficas}

PYTHON SOFTWARE FOUNDATION, 2024. \emph{hashlib --- Secure hashes and
message digests} {[}en línea{]}. Python 3.12 Documentation. Disponible
en: https://docs.python.org/3/library/hashlib.html {[}Consultado: 29 de
junio de 2026{]}.

VIRUSTOTAL / GOOGLE, 2024. \emph{YARA documentation} {[}en línea{]}.
GitHub Pages. Disponible en: https://virustotal.github.io/yara/
{[}Consultado: 29 de junio de 2026{]}.

FORTRA, 2024. \emph{impacket: A collection of Python classes for working
with network protocols} {[}en línea{]}. GitHub. Disponible en:
https://github.com/fortra/impacket {[}Consultado: 29 de junio de
2026{]}.

VAN ROSSUM, G. y DRAKE, F. L., 2009. \emph{Python 3 Reference Manual}
{[}en línea{]}. Python Software Foundation. Disponible en:
https://docs.python.org/3/reference/ {[}Consultado: 29 de junio de
2026{]}.

ÁLVAREZ, V. M., 2014. \emph{YARA: The pattern matching swiss knife for
malware researchers} {[}en línea{]}. VirusTotal Blog. Disponible en:
https://www.virustotal.com/gui/blog {[}Consultado: 29 de junio de
2026{]}.

SECUREAUTH CORPORATION, 2023. \emph{Impacket scripts and examples for
network protocol analysis} {[}en línea{]}. GitHub. Disponible en:
https://github.com/fortra/impacket/tree/master/examples {[}Consultado:
29 de junio de 2026{]}.

HUACAC SOTO, A., 2026. \emph{hashlib-Jara: Detección de APKs envenenados
con Python} {[}en línea{]}. GitHub. Disponible en:
https://github.com/AlvaroHuacacSoto/hashlib-Jara {[}Consultado: 29 de
junio de 2026{]}.

RAPID7, 2024. \emph{Metasploit Framework Documentation} {[}en línea{]}.
Disponible en: https://docs.metasploit.com/ {[}Consultado: 29 de junio
de 2026{]}.

ANDROGUARD, 2024. \emph{Androguard: Reverse engineering tool for Android
applications} {[}en línea{]}. GitHub. Disponible en:
https://github.com/androguard/androguard {[}Consultado: 29 de junio de
2026{]}.

\subsection{ANEXOS}\label{anexos}

\subsubsection{Anexo 1: Código fuente completo}\label{anexo-1-cuxf3digo-completo-del-notebook}

El código fuente completo de los scripts \texttt{comparar\_hashes.py} y
\texttt{detector\_yara.py} se encuentra en el repositorio GitHub:
\url{https://github.com/AlvaroHuacacSoto/hashlib-Jara}

\subsubsection{Anexo 2: Resultados de ejecución}\label{anexo-2-capturas-de-pantalla-de-la-ejecuciuxf3n}

\textbf{comparar\_hashes.py}

\begin{verbatim}
CairosOriginal.apk
  MD5     : 3392a4d54f89fc7fe9c6618ff418bd75
  SHA-1   : 1e4dc3d80c74e4a5f2aafb0cbd2de80119180708
  SHA-256 : 6752e6831ed260a7a53d79c4607084b392f31c6b0d527a446faacdbb283c2b7a
  SHA-512 : 201a5bacda49459a19fc3a66356ab3707d30d12ad1c4c7969c9b4415362fccd5...

CairosVenom.apk
  MD5     : 7a35a20a63573b0225c49be3a3f55295
  SHA-1   : cfb3bba3068982c74c180b3a980166a072938b4f
  SHA-256 : fe1452b4900c9214af83e81b7bada97a598a900be32fd9962f3a59c26c69abe4
  SHA-512 : 3708cf576beda849566697649088ecee4515d36b572239261bc447a99164ae25...

RESULTADO: Todos los hashes DISTINTOS.
Efecto avalancha SHA-256: 130/256 bits (50.8%)
\end{verbatim}

\textbf{detector\_yara.py}

\begin{verbatim}
CairosOriginal.apk: LIMPIO
CairosVenom.apk:    ENVENENADO
payload.apk:        ENVENENADO

APKs ENVENENADOS (2):
  - CairosVenom.apk (14 permisos peligrosos, subpaquete dfwto)
  - payload.apk     (14 permisos peligrosos)
APKs LIMPIOS (1):
  - CairosOriginal.apk (2 permisos)
\end{verbatim}

\subsubsection{Anexo 3: Archivos de prueba
utilizados}\label{anexo-3-archivos-de-prueba-utilizados}

\begin{itemize}
\tightlist
\item
  \texttt{APKs/CairosOriginal.apk} --- APK legítimo (52,400,958 bytes)
\item
  \texttt{APKs/CairosVenom.apk} --- APK con payload Metasploit
  (52,431,637 bytes)
\item
  \texttt{APKs/payload.apk} --- APK standalone de MsfVenom (10,246 bytes)
\end{itemize}

\subsubsection{Anexo 4: Enlaces a notebooks de Google
Colab}\label{anexo-4-enlaces-a-notebooks-de-google-colab}

\begin{itemize}
\tightlist
\item
  Notebook hashlib:\\
  \url{https://colab.research.google.com/drive/1oXyDNrLdfR4-8oltR1WURTp8IMK65TTF}
\item
  Notebook yara-python:\\
  \url{https://colab.research.google.com/drive/173T7NiHw52sU_ySdK0RQoJzm5KdplAmp}
\end{itemize}

\subsubsection{Anexo 5: Repositorio del
proyecto}\label{anexo-5-repositorio-del-proyecto}

\url{https://github.com/AlvaroHuacacSoto/hashlib-Jara}

Contiene: scripts Python, APKs de prueba, documento PDF y entorno
virtual con dependencias.

\subsubsection{Anexo 6: Regla YARA
utilizada}\label{anexo-6-regla-yara-utilizada}

\begin{verbatim}
rule dex_metasploit_payload {
    meta:
        description = "Detecta payload de Metasploit en DEX"
        author = "Lab Seguridad"
        severity = "critical"
    strings:
        $a = "Ljava/lang/reflect/Method;"
        $b = "Ljava/net/URLConnection;"
        $c = "addRequestProperty"
        $d = ";->startService(Landroid/content/Context;)V"
        $e = ";->openConnection()Ljava/net/URLConnection;"
    condition:
        uint32(0) == 0x0A786564 and
        4 of them
}

rule dex_payload_backdoor {
    meta:
        description = "Detecta estructura tipica de RAT/backdoor Android"
        severity = "high"
    strings:
        $ref = "Ljava/lang/reflect/Method;"
        $start = ";->startService(Landroid/content/Context;)V"
    condition:
        uint32(0) == 0x0A786564 and
        all of them
}
\end{verbatim}
\end{document}
